\documentclass[12pt,a4paper]{article}
 
% ── Paquetes ────────────────────────────────────────────────────────────────
\usepackage[utf8]{inputenc}
\usepackage[spanish]{babel}
\usepackage{amsmath, amssymb, amsthm}
\usepackage{booktabs}
\usepackage{array}
\usepackage{geometry}
\usepackage{graphicx}
\usepackage{xcolor}
\usepackage{float}
\usepackage{hyperref}
\usepackage{fancyhdr}
\usepackage{titlesec}
 
\geometry{margin=2.5cm}
 
% ── Colores ─────────────────────────────────────────────────────────────────
\definecolor{azuluni}{RGB}{0,70,127}
\definecolor{grisclaro}{RGB}{245,245,245}
 
% ── Encabezado ──────────────────────────────────────────────────────────────
\pagestyle{fancy}
\fancyhf{}
\rhead{\textcolor{azuluni}{\small T3 – Probabilidad Conjunta}}
\lhead{\textcolor{azuluni}{\small Caso Discreto}}
\rfoot{\thepage}
 
\titleformat{\section}{\large\bfseries\color{azuluni}}{}{0em}{}[\titlerule]
\titleformat{\subsection}{\normalsize\bfseries\color{azuluni}}{}{0em}{}
 
% ────────────────────────────────────────────────────────────────────────────
\begin{document}
 
% ── Portada ─────────────────────────────────────────────────────────────────
\begin{titlepage}
  \centering
  \vspace*{2cm}
  {\color{azuluni}\rule{\linewidth}{1.5pt}}\\[0.5cm]
  {\LARGE\bfseries\color{azuluni}
   Aplicación 1: Función de Probabilidad Conjunta\\[0.3cm]
   Caso Discreto}\\[0.3cm]
  {\color{azuluni}\rule{\linewidth}{1.5pt}}\\[1.5cm]
 
  {\large\bfseries Tarea T3 – Probabilidad Conjunta}\\[0.4cm]
  {\large Alejandro Escobar Barrios --- 20251020094
}\\[2cm]
 
  \begin{tabular}{ll}
    \textbf{Contexto:} & Defectos en producción industrial \\[0.2cm]
    \textbf{Variables:} & $X$: defectos en turno mañana \quad $Y$: defectos en turno tarde \\[0.2cm]
    \textbf{Herramientas:} & \LaTeX{} + Python (\texttt{numpy}, \texttt{matplotlib}) \\
  \end{tabular}
 
  \vfill
  {\small\color{gray} Generado con \LaTeX{} y validado con Python}
\end{titlepage}
 
% ────────────────────────────────────────────────────────────────────────────
\section{Planteamiento del Problema}
 
Una empresa manufacturera registra, en cada jornada laboral, el número de
piezas defectuosas producidas en el \textbf{turno mañana} ($X$) y en el
\textbf{turno tarde} ($Y$).  Ambas variables toman valores en
$\{0, 1, 2\}$ (0 = ninguna, 1 = una, 2 = dos o más).
 
A partir de datos históricos se construye la siguiente tabla de
\textbf{distribución de probabilidad conjunta} $P(X=x,\,Y=y)$:
 
\begin{table}[H]
\centering
\caption{Distribución de probabilidad conjunta $P(X=x,\,Y=y)$}
\renewcommand{\arraystretch}{1.4}
\begin{tabular}{c|ccc|c}
\toprule
\diagbox{$X$}{$Y$} & $y=0$ & $y=1$ & $y=2$ & $P(X=x)$ \\
\midrule
$x=0$ & 0.30 & 0.10 & 0.05 & 0.45 \\
$x=1$ & 0.15 & 0.20 & 0.05 & 0.40 \\
$x=2$ & 0.05 & 0.05 & 0.05 & 0.15 \\
\midrule
$P(Y=y)$ & 0.50 & 0.35 & 0.15 & \textbf{1.00} \\
\bottomrule
\end{tabular}
\end{table}
 
% ────────────────────────────────────────────────────────────────────────────
\section{Marco Teórico}
 
\subsection{Definición – Distribución Conjunta Discreta}
 
Sean $X$ e $Y$ variables aleatorias discretas definidas en el mismo espacio
muestral $\Omega$.  La \textbf{función de probabilidad conjunta} se define
como:
\[
  p(x,y) = P(X = x,\; Y = y), \qquad (x,y)\in\mathbb{Z}^{2}.
\]
 
Para que $p$ sea válida debe cumplir:
\begin{enumerate}
  \item $p(x,y)\geq 0$ para todo $(x,y)$.
  \item $\displaystyle\sum_{x}\sum_{y} p(x,y) = 1$.
\end{enumerate}
 
\subsection{Distribuciones Marginales}
 
Las distribuciones marginales se obtienen sumando sobre la otra variable:
\[
  p_X(x) = \sum_{y} p(x,y), \qquad
  p_Y(y) = \sum_{x} p(x,y).
\]
 
\subsection{Independencia}
 
$X$ e $Y$ son \textbf{independientes} si y solo si:
\[
  p(x,y) = p_X(x)\cdot p_Y(y) \quad \forall\,(x,y).
\]
 
\subsection{Esperanza, Varianza y Covarianza}
 
\[
  E[X] = \sum_{x} x\,p_X(x), \quad
  \text{Var}(X) = E[X^2] - (E[X])^2.
\]
\[
  \text{Cov}(X,Y) = E[XY] - E[X]\,E[Y], \qquad
  E[XY] = \sum_{x}\sum_{y} xy\,p(x,y).
\]
\[
  \rho_{XY} = \frac{\text{Cov}(X,Y)}{\sqrt{\text{Var}(X)\,\text{Var}(Y)}}.
\]
 
% ────────────────────────────────────────────────────────────────────────────
\section{Desarrollo Analítico}
 
\subsection{Verificación de la Distribución Conjunta}
 
\[
  \sum_{x=0}^{2}\sum_{y=0}^{2} p(x,y)
  = 0.30+0.10+0.05+0.15+0.20+0.05+0.05+0.05+0.05 = \mathbf{1.00} \checkmark
\]
 
\subsection{Distribuciones Marginales}
 
\textbf{Marginal de $X$:}
\[
  p_X(0)=0.45,\quad p_X(1)=0.40,\quad p_X(2)=0.15.
\]
 
\textbf{Marginal de $Y$:}
\[
  p_Y(0)=0.50,\quad p_Y(1)=0.35,\quad p_Y(2)=0.15.
\]
 
\subsection{Verificación de Independencia}
 
Se verifica si $p(0,0) = p_X(0)\cdot p_Y(0)$:
\[
  p_X(0)\cdot p_Y(0) = 0.45 \times 0.50 = 0.225 \neq 0.30 = p(0,0).
\]
$\therefore$ $X$ e $Y$ \textbf{no son independientes}.
 
\subsection{Esperanzas y Varianzas}
 
\[
  E[X] = 0(0.45)+1(0.40)+2(0.15) = 0.70
\]
\[
  E[X^2] = 0(0.45)+1(0.40)+4(0.15) = 1.00 \quad\Rightarrow\quad
  \text{Var}(X) = 1.00 - 0.49 = 0.51
\]
\[
  E[Y] = 0(0.50)+1(0.35)+2(0.15) = 0.65
\]
\[
  E[Y^2] = 0+0.35+0.60 = 0.95 \quad\Rightarrow\quad
  \text{Var}(Y) = 0.95 - 0.4225 = 0.5275
\]
 
\subsection{Covarianza y Correlación}
 
\[
  E[XY] = \sum_{x,y} xy\,p(x,y)
  = 1\!\cdot\!1(0.20)+1\!\cdot\!2(0.05)+2\!\cdot\!1(0.05)+2\!\cdot\!2(0.05)
  = 0.20+0.10+0.10+0.20 = 0.60
\]
\[
  \text{Cov}(X,Y) = 0.60 - (0.70)(0.65) = 0.60 - 0.455 = \mathbf{0.145}
\]
\[
  \rho_{XY} = \frac{0.145}{\sqrt{0.51\times 0.5275}}
            = \frac{0.145}{\sqrt{0.2690}}
            = \frac{0.145}{0.5187}
            \approx \mathbf{0.2796}
\]
 
La correlación positiva moderada indica que cuando aumentan los defectos en
el turno mañana, tienden a aumentar también en el turno tarde.
 
% ────────────────────────────────────────────────────────────────────────────
\section{Probabilidades Condicionales}
 
La distribución condicional de $Y$ dado $X=1$ es:
\[
  P(Y=y \mid X=1) = \frac{p(1,y)}{p_X(1)}
\]
\[
  P(Y=0\mid X=1)=\frac{0.15}{0.40}=0.375, \quad
  P(Y=1\mid X=1)=\frac{0.20}{0.40}=0.500, \quad
  P(Y=2\mid X=1)=\frac{0.05}{0.40}=0.125.
\]
 
La esperanza condicional:
\[
  E[Y\mid X=1] = 0(0.375)+1(0.500)+2(0.125) = 0.750.
\]
 
% ────────────────────────────────────────────────────────────────────────────
\section{Validación con Python}
 
El script \texttt{T3\_App1\_Discreta.py} (adjunto) reproduce todos los
cálculos anteriores y genera:
 
\begin{itemize}
  \item \textbf{Figura 1:} Gráfico de barras 3D de la distribución conjunta.
  \item \textbf{Figura 2:} Distribuciones marginales comparadas.
  \item \textbf{Figura 3:} Distribución condicional $P(Y\mid X=1)$.
\end{itemize}
 
Los resultados numéricos del script coinciden exactamente con los valores
obtenidos de forma analítica en las secciones anteriores.
 
% ────────────────────────────────────────────────────────────────────────────
\section{Conclusiones}
 
\begin{enumerate}
  \item La tabla de probabilidad conjunta es válida: suma 1 y todas las
        entradas son no negativas.
  \item Las variables $X$ e $Y$ \textbf{no son independientes}, lo que
        refleja que los procesos de los dos turnos están relacionados
        (e.g., misma maquinaria, mismos insumos).
  \item La covarianza positiva ($\text{Cov}=0.145$) y la correlación
        $\rho\approx 0.28$ cuantifican esta dependencia.
  \item La distribución condicional permite tomar decisiones: si el turno
        mañana produce 1 defecto, la probabilidad de que el turno tarde
        produzca 1 defecto sube al 50\%.
\end{enumerate}
 
\end{document}
